% 3SUM for engineers — practitioner version, plain language
% Build: tectonic paper.tex
\documentclass[11pt]{article}
\usepackage[margin=1.1in]{geometry}
\usepackage{amsmath}
\usepackage{booktabs}
\usepackage[hidelinks]{hyperref}
\usepackage{microtype}
\usepackage{listings}
\usepackage{xcolor}
\lstset{basicstyle=\ttfamily\small, columns=fullflexible, keepspaces=true,
        frame=single, framerule=0.3pt, xleftmargin=1em}

\title{Counting Triples That Sum to a Target:\\An Engineer's Field Guide to 3SUM}
\author{Melih Birim\\\texttt{melihbirim@gmail.com}}
\date{July 2026}

\begin{document}
\maketitle

\begin{abstract}
You have an array of integers and want to count how many triples add up to a
target, say zero. This is the classic 3SUM problem. We built four working
solutions in C, benchmarked them on the standard textbook inputs, and this paper
tells you which one to use and why, in plain language. The short version: sort the
array and use two pointers; that wins almost every time. But if your values live
in a limited range (say, sensor readings, prices in cents, pixel values) and you
have millions of them, there is a trick borrowed from signal processing whose
cost tracks the value range instead of the input size: on our inputs, about half
a second whether given one thousand numbers or sixty-four thousand, and 3 seconds
at ten million. We explain that trick the way one engineer would explain it to another,
show exactly where it starts winning (around 20{,}000--25{,}000 elements on our
inputs), and describe how we made sure every number in this paper is real: four
independent programs must agree with each other and with published reference
counts, and the whole benchmark re-runs in CI on every release. Code:
\url{https://github.com/melihbirim/nsum}.
\end{abstract}

\section{The problem, and why you should care}

Count the ways to pick three different positions $i<j<k$ in an array so that
$a_i+a_j+a_k$ equals a target $T$. It looks like a toy interview question. It is
also a real building block: collision detection, finding zero-net groups of
transactions, and dozens of computational geometry problems secretly reduce to it.

Theory says something discouraging: nobody knows how to solve general 3SUM
meaningfully faster than $n^2$ steps, and the standing conjecture is that nobody
ever will~\cite{gajentaan1995}. (The best known result shaves only
polylogarithmic factors off $n^2$~\cite{gronlund2014}; for an engineer, that is
quadratic.) So for arbitrary integers, quadratic is the game, and the engineering
question becomes: \emph{which} quadratic, and is there any way around it for
\emph{your} data? There is one way around it, and it only needs one property your
data may already have: a bounded value range.

\section{Four ways to do it}

\paragraph{1. Brute force: three nested loops.}
Try every triple. $O(n^3)$. On our machine it takes 16 seconds at 8{,}000 elements
and extrapolates to 17 minutes at 32{,}000 and over two hours at 64{,}000. Useless in production, but priceless
as ground truth: it is so simple it is almost certainly correct, so we use it to
check the clever versions on small inputs.

\paragraph{2. Hash map.}
For each pair, ask a hash table how many earlier elements equal $T - a_i - a_j$.
Still $O(n^2)$, and it needs no sorting. That is its only advantage. In our
measurements it is about $3\times$ slower than the next method, because hashing
jumps around memory while the alternative streams through it. Use it only when
you genuinely cannot sort (for example, data arriving as a stream you index once).

\paragraph{3. Sort + two pointers: the default.}
Sort once. For each element $a_k$, sweep two pointers toward each other across the
elements before it: sum too small, move the left pointer right; too big, move the
right pointer left; equal, count and move both. Each sweep is linear, so the whole
thing is $O(n^2)$ but with a tiny constant and perfectly sequential memory access.
One practical wrinkle: duplicated values. When both pointers sit on runs of equal
values, multiply the run lengths instead of stepping one at a time; when the whole
remaining window is a single repeated value appearing $m$ times, add
$m(m-1)/2$ pairs. That keeps repeated data from degrading the sweep.
\textbf{This is the method to reach for first.} It won every benchmark below
$20$K elements and it can list the triples, not just count them.

\paragraph{4. The FFT trick: the escape hatch.}
The only way we know to beat quadratic, and it works only when values sit in a
bounded range $U$ (in our inputs, $[-10^6, 10^6]$, so $U \approx 2$ million). Its
running time depends on $U$, \emph{not on $n$}: about half a second at every size
from 1K to 64K on our machine. (At 10 million elements the linear work of reading
input and building the histogram finally shows up, and the run takes about 3
seconds. Against 36 extrapolated hours for the alternative, we will take it.)
Section~\ref{sec:fft} explains the trick without assuming you remember any
signal-processing math. The trick itself is well known, folklore in competitive
programming and a textbook exercise; our contribution is measuring it against the
alternatives in production C and spelling out the exact-counting correction.

\paragraph{A detour that failed: SIMD.}
We first tried to speed up brute force with NEON vector instructions, comparing
four integers per instruction. It genuinely made the inner loop $4\times$ faster.
It was still hopeless, because two-pointer's \emph{algorithmic} advantage is a
factor around $n/6$, which at $n=8$K is over a thousand. We kept the code in the
repo as a reminder of the rule every performance engineer eventually learns the
hard way: \textbf{pick the right algorithm first, then vectorize it.} A $4\times$
constant does not rescue the wrong big-O.

\section{The FFT trick, explained like an engineer}
\label{sec:fft}

Forget the array order. Build a histogram: \texttt{cnt[v]} = how many times value
$v$ appears. All the information 3SUM counting needs is in that histogram.

Now the key idea. Suppose you want a table \texttt{pairs[s]} = number of ways to
pick two values (with repetition, order counted) that add up to $s$. You could
compute it with two nested loops over the histogram, which is $O(U^2)$, quadratic
again, just in $U$ instead of $n$. But this exact table is what mathematicians
call a \emph{convolution} of the histogram with itself, and there is a standard
fast routine for convolutions: the Fast Fourier Transform. You do not need to
understand \emph{why} the FFT works to use it, the same way you do not need to
understand how your compiler allocates registers. The contract is:

\begin{lstlisting}
FFT(histogram)          // transform: O(U log U)
square it pointwise     // O(U)
inverse FFT             // back-transform: O(U log U)
result = pairs[]        // pair-sum table, all sums at once
\end{lstlisting}

That is the whole trick: \textbf{the FFT is a fast way to fill in the ``how many
pairs sum to $s$'' table for every $s$ simultaneously.} With that table, counting
triples is one more linear pass: for each value $c$ in the histogram, look up
\texttt{pairs[T - c]} and multiply by \texttt{cnt[c]}.

\paragraph{The correction step.}
There is a catch, and it is the same catch you hit whenever you count
combinations with a tool that does not know about ``distinct positions.'' The
pair table happily counts an element paired \emph{with itself}, and our triple
loop happily reuses the same position twice. So the raw total, call it $S_1$,
over-counts. The fix is ordinary double-counting bookkeeping, no advanced math:

\begin{itemize}
\item Count triples that reuse one position twice: for each value $c$, that needs
      a partner equal to $T - 2c$; call the total $A$.
\item Count triples that use one position three times: that is only possible if
      $T$ is divisible by 3, and equals \texttt{cnt[T/3]}; call it $B$.
\item Every over-count pattern can happen in a few symmetric ways, and every
      legitimate triple got counted $6$ times ($3!$ orderings). Putting it
      together:
\end{itemize}
\[
\text{answer} \;=\; \frac{S_1 - 3A + 2B}{6}.
\]
If you have ever de-duplicated a JOIN that matched rows to themselves, you have
done this kind of correction before. The $A$ and $B$ passes are cheap linear
loops. Total cost: $O(n)$ to build the histogram plus $O(U \log U)$ for the FFT.

\paragraph{One honest caveat.}
Our FFT uses \texttt{double} floats and rounds each table entry to the nearest
integer. Through 64K elements the counts are verified exactly right against the
integer-arithmetic methods. At 10 million elements the count
(12{,}499{,}729{,}168{,}317) is beyond independent verification: the quadratic
methods would need a day and a half to check it. A back-of-envelope error bound
says the rounding margin is comfortable there, but honesty requires the label
``unverified.'' If you operate at that scale for real, switch the FFT for an
integer version (an NTT); the recipe above does not change, only the transform
underneath.

\section{Measurements: who wins, where}

Standard benchmark inputs from Sedgewick and Wayne's \emph{Algorithms}
textbook~\cite{sedgewick2011} (the \texttt{ThreeSum} files, values in
$[-10^6,10^6]$), counting zero-sum triples, single-threaded C at \texttt{-O3} on
Apple arm64.

\begin{table}[h]
\centering
\caption{Seconds to count all zero-sum triples. $*$ = extrapolated; brute force
is unrunnable past 8K.}
\begin{tabular}{rrrrr}
\toprule
$n$ & brute & hash & two-pointer & FFT\\
\midrule
1K  & 0.035        & 0.003  & 0.0013 & 0.45\\
8K  & 16.2         & 0.23   & 0.083  & 0.48\\
32K & $\sim$1040$^*$  & 4.3    & 1.4    & 0.49\\
64K & $\sim$8300$^*$  & 17.4   & 5.5    & 0.45\\
\bottomrule
\end{tabular}
\end{table}

Read the FFT column top to bottom: it barely moves. Its cost is pinned to the
value range, and the range is the same in every file. Meanwhile the two-pointer
column grows $4\times$ per doubling, exactly as $n^2$ predicts. The lines cross
between 8K and 32K, around $n \approx 20$--$25$K on this hardware. Past the
crossing the gap explodes: at 10 million bounded values, two-pointer extrapolates
to about 36 hours; the FFT counter, actually run, finished in 3.1 seconds.

\paragraph{The decision table.}

\begin{itemize}
\item Everyday case, need to list the triples, or $n$ under $\sim$20K:
      \textbf{sort + two pointers.}
\item Can't sort: \textbf{hash}, and accept roughly $3\times$ slower.
\item Millions of elements whose values fit a bounded range: \textbf{FFT.}
\item Huge value range (64-bit hashes, arbitrary integers): FFT is out; its
      memory and time follow the range. Back to two pointers.
\end{itemize}

\section{How we made sure these numbers are real}

Benchmarks are easy to get wrong and easier to transcribe wrong, so the
repository enforces three rules, and we recommend them for any performance
writeup:

\begin{enumerate}
\item \textbf{Independent implementations must agree.} All four methods run on
      the same input; any disagreement prints \texttt{MISMATCH} and fails the
      run. Four programs written differently agreeing on 16{,}377{,}796 is strong
      evidence there is no shared bug.
\item \textbf{Check against numbers we didn't produce.} The textbook publishes
      the triple counts for its own inputs (70, 528, 4039, 32074, 255181,
      2052358, 16377796). Ours match, all seven.
\item \textbf{Let a machine that isn't yours reproduce it.} The release CI runs
      the full benchmark on a clean runner and pastes the actual output into the
      release notes. Numbers in this paper were reproduced, not copied.
\end{enumerate}

Each binary also carries a tiny built-in self-test (\texttt{demo()}) asserting
known answers on hand-checkable arrays, including the all-equal case
$[5,5,5]$, target 15, which exercises the $B$ correction.

\section{Takeaways}

\begin{enumerate}
\item Sort + two pointers is the production default for 3SUM. Simple, fast,
      lists triples, handles duplicates with the block-counting trick.
\item Vectorizing the wrong algorithm is wasted effort. Algorithm first,
      intrinsics second.
\item The FFT is not just for signals. Treat it as a library routine that fills
      an entire ``how many pairs sum to $s$'' table in $O(U\log U)$, and a whole
      class of counting problems drops out of quadratic time, whenever your
      values live in a bounded range.
\item Publish benchmarks with their verification harness. Cross-validation,
      external reference counts, and CI re-execution cost an afternoon and make
      the numbers trustworthy.
\end{enumerate}

\begin{thebibliography}{9}

\bibitem{gajentaan1995}
A.~Gajentaan and M.~H. Overmars.
\newblock On a class of $O(n^2)$ problems in computational geometry.
\newblock \emph{Computational Geometry}, 5(3):165--185, 1995.

\bibitem{gronlund2014}
A.~Gr{\o}nlund and S.~Pettie.
\newblock Threesomes, degenerates, and love triangles.
\newblock In \emph{Proc.\ 55th IEEE FOCS}, pages 621--630, 2014.

\bibitem{sedgewick2011}
R.~Sedgewick and K.~Wayne.
\newblock \emph{Algorithms}, 4th edition.
\newblock Addison-Wesley, 2011. \texttt{ThreeSum} benchmark data,
  \url{https://algs4.cs.princeton.edu}.

\end{thebibliography}

\end{document}
